%=============================================================================!
% 8.1  WHERE THE POTENTIAL ENERGY COMES FROM                                  !
%=============================================================================!
\section{Where the potential energy comes from}
\label{sec:pe-origin}

An atom is a small, heavy nucleus surrounded by much lighter electrons.
The forces between atoms are the electrical attractions and repulsions of
these charges: nuclei repel nuclei, electrons repel electrons, and each
attracts the other. This section explains how the motion of the electrons
can be set aside, leaving a potential energy for the nuclei alone, and
what it costs to compute it.

\subsection*{Heavy nuclei, light electrons}

A proton is 1836 times as heavy as an electron, and a lithium atom, of
mass \SI{6.94}{amu} (\cref{sec:ne-atoms}), 12\,651 times, almost all of
it in the nucleus. Pushed by forces of similar size, a body's acceleration is
inversely proportional to its mass, and \cref{sec:os-small} showed that a
body held by a spring of stiffness $k$ swings with $\omega_0 =
\sqrt{k/m}$: with comparable stiffnesses, the electrons respond
$\sqrt{1836} = 43$ times faster than a proton and 112 times faster than
lithium. On the time scale on which the nuclei move, the electrons
have time to settle into the arrangement of lowest energy for wherever
the nuclei happen to be, and to follow them as they move.

This is the picture of Born and Oppenheimer\cite{BornOppenheimer1927}.
With the nuclei held still at the positions $\vb{r}^N$, the electrons
settle into their state of lowest energy for those positions, and that
energy, together with the repulsion of the nuclei for one another,
depends only on $\vb{r}^N$; we call it $U(\vb{r}^N)$. When the nuclei move, the electrons
follow, and the nuclei move as if pushed by the forces $-\grad_iU$
(\cref{sec:en-atoms}). The function $U$ is called the potential energy
surface: for a single coordinate it is a curve, for two a landscape like
that of the lithium atom on its model layer (\cref{sec:en-surface}), and
for $N$ atoms a surface in $3N$ dimensions. The picture fails when the
motion of the nuclei carries the electrons from one state into another
close in energy, as after light is absorbed or in some chemical
reactions; every system in this book stays in its lowest state.

\subsection*{The cost of computing it}

Finding the lowest state of the electrons is a problem of quantum
mechanics. Density-functional theory, one such method, turns it into the eigenvalues and eigenvectors (\cref{sec:os-modes}) of a
matrix whose size $n$ grows in proportion to the number of atoms. The
work of finding all of them grows as $n^3$, and \cref{fig:pe-cost}(a)
measures it:
with NumPy, doubling the size of a dense symmetric matrix multiplied the
time to find its eigenvalues and eigenvectors by about 7.5, close to $2^3 =
8$, and the times for the larger matrices lie near a line of slope 3 on
logarithmic axes
(\cref{sec:mo-taylor}): a time proportional to $n^3$ has a logarithm that
grows by 3 times the logarithm of $n$. The energy of a model that adds up the energies of
pairs of atoms needs one term for each pair: each of the $N$ atoms pairs
with $N - 1$ others, which counts every pair twice, so there are $N(N -
1)/2 \approx N^2/2$ pairs for large $N$, a number that doubling $N$
multiplies by $2^2 = 4$; doubling $N$ multiplied its time by about 4
(\cref{fig:pe-cost}b); with the cutoffs of \cref{sec:pe-cutoffs} and the
neighbour lists of Chapter~10 it grows only as $N$. Computing $U$ from the
electrons at every step of a simulation therefore limits it to far
smaller systems and shorter times than a model allows, and most
simulations use a model of $U$ fitted to such calculations, which is the subject of the rest of this
chapter.

\begin{figure}[htb]
\centering
\includegraphics{ch08_potentials/cost}
\caption{What computing the energy costs, with NumPy on one processor
core; the times depend on the computer, the slope at large sizes does
not. (a) Finding
all the eigenvalues and eigenvectors of a dense symmetric matrix of size
$n$ (one with no zero entries to exploit), against $n$, with a line of
slope 3 (dashed). (b) The energy and forces of $N$ atoms interacting in
pairs, every pair computed, against $N$, with a line of slope 2
(dashed).}
\label{fig:pe-cost}
\end{figure}

\begin{bridge}
The forces on the nuclei follow from the electronic energy by the
Hellmann--Feynman theorem\cite{Feynman1939}: at the converged density they
are the electrostatic forces of the electrons and the other nuclei, plus
Pulay corrections when the basis functions move with the atoms. The
cubic cost is that of a dense eigensolver; linear-scaling methods use the
nearsightedness of the density matrix to reach a cost proportional to
$N$\cite{Kohn1996}, and machine-learned potentials (planned for the
continuation of Part~V) build the same
locality into a model fitted to the expensive calculations.
\end{bridge}

\begin{keyidea}
Because the electrons are light and fast, they settle at every
arrangement of the nuclei, and the energy of that settled state, with the
repulsion of the nuclei, is the potential energy surface $U(\vb{r}^N)$ on
which the nuclei move. Computing it from the electrons costs, in the usual
methods, work growing as the cube of the system's size, which is why most simulations fit a
model of it.
\end{keyidea}

\begin{selfcheck}
Finding all the eigenvalues and eigenvectors of a matrix of size 1000
takes \SI{0.1}{\second}. Roughly how long does it take for one of size
4000?
\end{selfcheck}
